% en/lessons/02-lesson-skeleton.tex の日本語版．英語版と節・ラベル・
% 余白ノートを完全に対応させる．規則: docs/style-guide.md．
\documentclass[../main.tex]{subfiles}
\begin{document}

\chapter{章の雛形}
\lessoninfo{
  video={Lesson 2，動画 1--12},
  textbook={H\&P \cite{hennessy2017} 第1章},
  keywords={最初の用語，次の用語},
}

2〜4文で，この章で扱う内容，その意義，前の章とのつながりを述べる．

\section{最初のトピック}
\label{sec:l02-first}

節の最初の文．\source{Lesson 2，動画 1--5；H\&P §1.2．}
\term[さいしょのようご]{最初の用語}[first term]は，資料が定義する箇所で定義する．
説明は資料の順序に従い，概念が何であるか，なぜ必要か，どう使われるかを述べる．

\begin{definition}[最初の用語]
  \label{def:l02-first}
  定義の正確な記述．
\end{definition}

\begin{example}
  \label{ex:l02-first}
  独自の数値を用いた計算例．\margin{例で使う記号と単位．}
\end{example}

\section{次のトピック}
\label{sec:l02-second}

\cref{sec:l02-first}を踏まえた本文．\source{Lesson 2，動画 6--12．}
余白の項目は，それが補う文の隣に置き，短い文章の範囲には多くても3つ程度に
とどめる．

経験則はヒントにする．\tip{経験則．}余白を読み飛ばしても議論を追えるように，
論証は本文で続ける．3行を超える補足は番号付きの注にする．\sidenote{設計の
理由や，実在する実装の名前とその出典．}

\begin{keypoint}[章のまとめ]
  \begin{itemize}
    \item 最初の節で覚えておくべき一点．
    \item 次の節で覚えておくべき一点．
  \end{itemize}
\end{keypoint}

\end{document}
